% ------------------------------------------------------------------------------
\section{Solución de problemas}
% ------------------------------------------------------------------------------

\begin{longtable}{L{5.2cm}L{8.6cm}}
	\caption{Problemas frecuentes al instalar y abrir ConcretoLab.} \label{tab:problemas} \\
	\hline
	\textbf{Lo que ve} & \textbf{Qué hacer} \\ \hline
	\endfirsthead
	\hline
	\textbf{Lo que ve} & \textbf{Qué hacer} \\ \hline
	\endhead
	\hline
	\endfoot

	«Windows protegió su PC» y no encuentra cómo continuar. &
	Pulse el texto \boton{Más información}: recién entonces aparece \boton{Ejecutar de todas formas}
	(paso~2, página~\pageref{sec:smartscreen}). \\ \hline

	El antivirus bloquea o pone en cuarentena el instalador. &
	Ocurre a veces con programas sin firma digital. Restaure el archivo desde el antivirus o pida a soporte técnico que
	lo permita, y vuelva a ejecutarlo. Si duda del archivo, compruebe su huella (anexo~\ref{sec:huella}). \\ \hline

	«Otra instalación está en curso» o el asistente no avanza. &
	Espere a que termine la otra instalación (por ejemplo, actualizaciones de Windows) o reinicie la computadora y vuelva a
	empezar. \\ \hline

	La carpeta de destino ya existe. &
	El asistente pregunta si desea instalar ahí de todos modos: si es una instalación anterior de ConcretoLab, responda
	\boton{Sí}. \\ \hline

	No encuentra el programa después de instalar. &
	Pulse la tecla \tecla{Windows}, escriba \textbf{ConcretoLab} y pulse \tecla{Enter}. \\ \hline

	Al abrir, tarda en aparecer la ventana. &
	Es normal que la primera vez tarde unos segundos más. Si no aparece después de un minuto, reinicie la computadora y
	vuelva a intentarlo. \\ \hline

	Pide una contraseña al abrir. &
	Sus datos son de una versión anterior que usaba contraseña. Escríbala una sola vez (paso~3 de la sección «Abrir
	ConcretoLab por primera vez»). Si no la recuerda, pulse \boton{Salir}: sus datos no se modifican. \\ \hline

	«No se pudo convertir la base». &
	Sus datos no se modificaron. Cierre otros programas que puedan estar usando la carpeta \ruta{.concretolab} (por ejemplo,
	una copia en curso) y vuelva a abrir ConcretoLab. \\ \hline

	Quiere llevar sus proyectos a otra computadora. &
	Dentro del programa, use \menu{Datos y copias} \textgreater{} \boton{Exportar...} y en la otra computadora
	\boton{Importar...}. Se explica en el Manual de usuario. \\ \hline

	No tiene espacio suficiente en el disco. &
	Libere al menos 300~MB (por ejemplo, vaciando la papelera) y vuelva a ejecutar el instalador. \\
\end{longtable}

Si el problema continúa, anote el mensaje exacto que aparece en pantalla y comuníquese con el equipo de ConcretoLab. Los
registros técnicos de la carpeta \ruta{.concretolab\textbackslash logs} ayudan a encontrar la causa.
